\section{Kinematics}
\label{sec:kinematics}

Each cell is a translational parallelogram delta actuated by three prismatic
joints whose axes are vertical and mutually parallel. The parallelogram limbs
hold the moving platform parallel to the base, so the platform has three
translational degrees of freedom and no orientation state, and each limb reduces
to a single distance constraint.

\subsection{Notation}

Let the base frame be centred on the base-triangle centroid with
$\hat{\mathbf{z}}$ along the actuator axes, and index the limbs by azimuth
$\theta_i \in \{\pi/2,\, -\pi/6,\, 7\pi/6\}$, $\hat{\mathbf{u}}_i =
[\cos\theta_i,\ \sin\theta_i]^{\!\top}$. With base and platform circumradii
$R_b=s_b/\sqrt3$ and $R_p=s_p/\sqrt3$, prismatic coordinates
$\mathbf{q}=[q_1,q_2,q_3]^{\!\top}$ and platform centre
$\mathbf{x}=[x,y,z]^{\!\top}$, the carriage and platform joints of limb $i$ are
\begin{equation}
\mathbf{c}_i = R_b \begin{bmatrix}\hat{\mathbf{u}}_i\\0\end{bmatrix} + q_i\hat{\mathbf{z}},
\qquad
\mathbf{p}_i = \mathbf{x} + R_p \begin{bmatrix}\hat{\mathbf{u}}_i\\0\end{bmatrix}.
\end{equation}
The tool tip is offset from the platform plane by a constant $d$ along the
invariant platform normal, $\mathbf{x}_{\text{ee}}=\mathbf{x}+d\hat{\mathbf{z}}$.
For limbs of length $L$, loop closure is
\begin{equation}
\label{eq:closure}
\|\mathbf{p}_i-\mathbf{c}_i\| = L, \qquad i=1,2,3.
\end{equation}

\subsection{Inverse Kinematics}

Base anchor and platform joint share the azimuth $\theta_i$, so their difference
collapses onto $\hat{\mathbf{u}}_i$ and the horizontal part of
\eqref{eq:closure} depends only on $R_p-R_b$:
\begin{equation}
\label{eq:rho}
\rho_i^2 = \big\|\,[x,\,y]^{\!\top} + (R_p-R_b)\,\hat{\mathbf{u}}_i \,\big\|^2 .
\end{equation}
Equation \eqref{eq:closure} then separates into three scalar equations
$\rho_i^2+(z-q_i)^2=L^2$, each in one actuator, giving the closed-form,
fully decoupled inverse map
\begin{equation}
\label{eq:ik}
q_i = \big(z_{\text{ee}}-d\big) - \sqrt{L^2-\rho_i^2}.
\end{equation}
The negative branch places each carriage below its platform joint, which is the
assembly mode of the prototype and renders the solution unique. The reachable
set is $\{\mathbf{x} : \rho_i \le L,\ q_i \in [q_{\min},q_{\max}]\ \forall i\}$.

\subsection{Forward Kinematics}

Absorbing the platform offset into the carriage position,
$\mathbf{s}_i = \mathbf{c}_i - R_p[\hat{\mathbf{u}}_i^{\!\top},0]^{\!\top}$,
turns \eqref{eq:closure} into $\|\mathbf{x}-\mathbf{s}_i\|=L$: the platform
centre lies at the intersection of three spheres of equal radius, centred on a
vertical prism of circumradius $R_b-R_p$ at heights $q_i$. Let
$\mathbf{v}_2=\mathbf{s}_2-\mathbf{s}_1$, $\mathbf{v}_3=\mathbf{s}_3-\mathbf{s}_1$,
and build the orthonormal frame
\begin{equation}
\hat{\mathbf{e}}_1=\frac{\mathbf{v}_2}{D},\quad
\hat{\mathbf{e}}_2=\frac{\mathbf{v}_3-\alpha\hat{\mathbf{e}}_1}{\beta},\quad
\hat{\mathbf{e}}_3=\hat{\mathbf{e}}_1\times\hat{\mathbf{e}}_2,
\end{equation}
with $D=\|\mathbf{v}_2\|$, $\alpha=\mathbf{v}_3\!\cdot\!\hat{\mathbf{e}}_1$ and
$\beta=\|\mathbf{v}_3-\alpha\hat{\mathbf{e}}_1\|$. Because the three radii are
equal, the pairwise sphere differences are linear and solve immediately:
\begin{equation}
\label{eq:fk}
\mathbf{x}=\mathbf{s}_1+X\hat{\mathbf{e}}_1+Y\hat{\mathbf{e}}_2+Z\hat{\mathbf{e}}_3,
\end{equation}
\begin{equation}
X=\frac{D}{2},\quad
Y=\frac{\|\mathbf{v}_3\|^2-\alpha D}{2\beta},\quad
Z=\sigma\sqrt{L^{2}-X^{2}-Y^{2}},
\end{equation}
where $\sigma=\operatorname{sgn}(\hat{\mathbf{e}}_3^{\!\top}\hat{\mathbf{z}})$
selects the branch with the platform above the plane of the joint centres,
consistent with \eqref{eq:ik}. The map is defined whenever $\beta\neq0$ and
$X^2+Y^2\le L^2$. The first condition is structural: $D\beta =
\|\mathbf{v}_2\times\mathbf{v}_3\| \ge 2A$, where $A$ is the constant area of
the equilateral triangle onto which the centres project, so collinearity cannot
occur. The second delimits the set of actuator triples admitting a rigid
platform.

\subsection{Jacobian and Singularities}

Differentiating \eqref{eq:closure} with limb vectors
$\mathbf{w}_i=\mathbf{p}_i-\mathbf{c}_i$ gives
$\mathbf{w}_i^{\!\top}\dot{\mathbf{x}}=w_{i,z}\dot{q}_i$, where
$w_{i,z}=\sqrt{L^2-\rho_i^2}$, hence
$\mathbf{J}_q\dot{\mathbf{q}}=\mathbf{J}_x\dot{\mathbf{x}}$ with
\begin{equation}
\label{eq:jac}
\mathbf{J}_q=\operatorname{diag}(w_{i,z}),\qquad
\mathbf{J}_x=\big[\mathbf{w}_1\ \mathbf{w}_2\ \mathbf{w}_3\big]^{\!\top}.
\end{equation}
Inverse-kinematic singularities ($\det\mathbf{J}_q=0$) require a horizontal limb,
$\rho_i\to L$; direct-kinematic singularities ($\det\mathbf{J}_x=0$) require the
three limb vectors to become linearly dependent. Sampled over the reachable set,
$\min_i w_{i,z}\ge1.78$\,mm and $|\det\mathbf{J}_x|\ge4.3\times10^{-6}$\,m$^3$,
so the workspace is singularity-free.

At the home pose $\mathbf{J}=\mathbf{J}_x^{-1}\mathbf{J}_q$ has singular values
$(8.39,\,8.39,\,0.577)$ and condition number $14.5$, rising to $135$ at the
workspace boundary. Vertical motion is the unity-gain common mode,
$\partial z/\partial q_i=1/3$, whereas differential actuation is amplified
radially by $\approx\!6.8$. Actuator resolution and error therefore appear in
the plane roughly an order of magnitude larger than along $\hat{\mathbf{z}}$,
which bounds lateral placement accuracy.

\begin{table}[t]
\caption{Kinematic parameters of one delta.}
\label{tab:params}
\centering
\begin{tabular}{llr}
\toprule
Symbol & Quantity & Value (mm) \\
\midrule
$s_b$ & base triangle side & 21.39 \\
$s_p$ & platform triangle side & 12.00 \\
$L$   & limb length & 56.00 \\
$d$   & platform plane to tool tip & 24.00 \\
$q_{\min},q_{\max}$ & prismatic travel & 5.0,\ 98.5 \\
\bottomrule
\end{tabular}
\end{table}

The limbs attach at the midpoints of the end-effector triangle's sides, so by
the midsegment theorem $s_p$ is half the end-effector side length. Home is taken
at mid-travel, placing the tip at $(0,0,131.49)$\,mm; the workspace spans
$z\in[85,178]$\,mm on the central axis, with a maximum radius of $53$\,mm.
